\Article{18}
\textsc{Aufgabe.} \textit{Es bezeichne \(U\) eine gegebene Funktion der unbekannten Größen \(V,V',V''\) etc.; gesucht ist der mittlere Fehler \(M\), der bei der Bestimmung des Wertes von \(U\) zu befürchten ist, wenn für \(V,V',V''\) etc. nicht die wahren Werte, sondern diejenigen angenommen werden, die aus voneinander unabhängigen Beobachtungen hervorgehen, deren mittlere Fehler jeweils \(m,m',m''\) etc. sind.}

\textit{Lösung.} Bezeichnen wir die Fehler in den beobachteten Werten von \(V,V',V''\) etc. mit \(e,e',e''\) etc., so lässt sich der daraus im Wert von \(U\) entstehende Fehler durch die lineare Funktion
\[
\lambda e+\lambda'e'+\lambda''e''+\text{etc.}=E
\]
ausdrücken, wobei \(\lambda,\lambda',\lambda''\) etc. die Werte der Differentialquotienten \linebreak[4]\(\frac{\mathrm{d}U}{\mathrm{d}V},\frac{\mathrm{d}U}{\mathrm{d}V'},\frac{\mathrm{d}U}{\mathrm{d}V''}\) etc. für die wahren Werte von \(V,V',V''\) etc. sind, sofern die Beobachtungen genau genug sind, dass die Quadrate und Produkte der Fehler vernachlässigt werden dürfen. Da die Beobachtungsfehler als frei von konstanten Teilen vorausgesetzt werden, folgt hieraus zunächst, dass der Mittelwert von \(E\) \(=0\) ist. Ferner wird der im Wert von \(U\) zu befürchtende mittlere Fehler die Quadratwurzel aus dem Mittelwert von \(EE\) sein, oder \(MM\) wird der Mittelwert der Summe
\[
\begin{aligned}
\lambda\lambda ee+\lambda'\lambda'e'e'+\lambda''\lambda''e''e''+\text{etc.}{}\\
{}+2\lambda\lambda'ee'+2\lambda\lambda''ee''+2\lambda'\lambda''e'e''+\text{etc.}
\end{aligned}
\]
sein. Der Mittelwert von \(\lambda\lambda ee\) wird aber \(\lambda\lambda mm\), der Mittelwert von \(\lambda'\lambda'e'e'\) wird \(=\lambda'\lambda'm'm'\) etc.; schließlich werden die Mittelwerte der Produkte \(2\lambda\lambda'ee'\) etc. sämtlich \(=0\). Hieraus schließen wir also
\[
M=\BareRoot{(\lambda\lambda mm+\lambda'\lambda'm'm'+\lambda''\lambda''m''m''+\text{etc.})}
\]
\SourceLeaf{025}\par
Es empfiehlt sich, dieser Lösung einige Bemerkungen hinzuzufügen.

I. Soweit die Beobachtungsfehler als Größen erster Ordnung betrachtet und Größen höherer Ordnungen vernachlässigt werden, dürfen in unserer Formel für \(\lambda,\lambda',\lambda''\) etc. auch diejenigen Werte der Quotienten \(\frac{\mathrm{d}U}{\mathrm{d}V}\) etc. angenommen werden, die sich aus den beobachteten Werten der Größen \(V,V',V''\) etc. ergeben. Ist \(U\) eine lineare Funktion, so wird offensichtlich überhaupt kein Unterschied bestehen.

II. Wollen wir statt der mittleren Fehler der Beobachtungen lieber deren Gewichte einführen, so seien diese, nach einer beliebigen Einheit, jeweils \(p,p',p''\) etc.; ferner sei \(P\) das Gewicht der Bestimmung des Wertes von \(U\) aus den beobachteten Werten der Größen \(V,V',V''\) etc. So erhalten wir
\[
P=\frac{1}{\dfrac{\lambda\lambda}{p}+\dfrac{\lambda'\lambda'}{p'}+\dfrac{\lambda''\lambda''}{p''}+\text{etc.}}
\]

III. Ist \(T\) eine andere gegebene Funktion der Größen \(V,V',V''\) etc. und gilt für deren wahre Werte
\[
\frac{\mathrm{d}T}{\mathrm{d}V}=\kappa,\quad\frac{\mathrm{d}T}{\mathrm{d}V'}=\kappa',\quad\frac{\mathrm{d}T}{\mathrm{d}V''}=\kappa''\text{ etc.}
\]
so wird der Fehler bei der Bestimmung des Wertes von \(T\) aus den beobachteten Werten von \(V,V',V''\) etc. \(=\kappa e+\kappa'e'+\kappa''e''+\text{etc.},\ =E'\) sein und der dabei zu befürchtende mittlere Fehler \(=\BareRoot{(\kappa\kappa mm+\kappa'\kappa'm'm'+\kappa''\kappa''m''m''+\text{etc.})}\). Die Fehler \(E,E'\) werden aber offensichtlich nicht mehr voneinander unabhängig sein. Der betreffende Mittelwert ist der des Produkts \(EE'\). Anders als der Mittelwert des Produkts \(ee'\) wird er nicht \(=0\), sondern \(=\kappa\lambda mm+\kappa'\lambda'm'm'+\kappa''\lambda''m''m''+\text{etc.}\) sein.

IV. Unsere Aufgabe lässt sich auch auf den Fall ausdehnen, in dem die Werte der Größen \(V,V',V''\) etc. nicht unmittelbar durch Beobachtungen gefunden, sondern auf beliebige Weise aus Kombinationen von Beobachtungen abgeleitet werden, sofern nur die einzelnen Bestimmungen voneinander unabhängig sind, das heißt auf verschiedenen Beobachtungen beruhen: Wenn diese Bedingung nicht erfüllt ist, wäre die Formel für \(M\) fehlerhaft. Wäre zum Beispiel die eine oder andere Beobachtung, die zur Bestimmung des Wertes von \(V\) gedient hat, auch zur Bestimmung des Wertes von \(V'\) verwendet worden, so wären die Fehler \(e\) und \(e'\) nicht mehr voneinander unabhängig und somit auch der Mittelwert des Produkts \(ee'\) nicht \(=0\). Wird aber in einem solchen Fall der Zusammenhang der Größen \(V,V'\) mit den einzelnen Beobachtungen, aus denen sie abgeleitet sind, gehörig berücksichtigt, so lässt sich der \SourceLeaf{026}Mittelwert des Produkts \(ee'\) mithilfe von Bemerkung III. angeben und damit die Formel für \(M\) vervollständigen.

\Article{19}
Es seien \(V,V',V''\) etc. Funktionen der Unbekannten \(x,y,z\) etc., die Anzahl der ersteren \(=\pi\), die Anzahl der Unbekannten \(=\rho\); wir nehmen an, dass durch unmittelbare oder mittelbare Beobachtungen die Werte der Funktionen als \(V=L,V'=L',V''=L''\) etc. gefunden worden sind, wobei diese Bestimmungen jedoch voneinander unabhängig waren. Ist \(\rho\) größer als \(\pi\), so ist die Ermittlung der Unbekannten offensichtlich eine unbestimmte Aufgabe; der nächste Fall liegt vor, wenn \(\rho\) gleich \(\pi\) ist. Jede der Unbekannten \(x,y,z\) etc. kann auf die Form einer Funktion von \(V,V',V''\) etc. gebracht oder als darauf gebracht gedacht werden, sodass sich ihre Werte aus den beobachteten Werten der letzteren finden lassen; zugleich kann mithilfe des vorigen Artikels die relative Genauigkeit angegeben werden, die jeder dieser Bestimmungen zukommt. Ist schließlich \(\rho\) kleiner als \(\pi\), so lässt sich jede der Unbekannten \(x,y,z\) etc. auf unendlich viele verschiedene Arten in die Form einer Funktion von \(V,V',V''\) etc. bringen und ihr Wert somit auf unendlich viele verschiedene Arten ermitteln. Diese Bestimmungen müssten freilich genau übereinstimmen, wenn die Beobachtungen absolute Genauigkeit hätten; da dies nicht der Fall ist, werden verschiedene Verfahren verschiedene Werte liefern, und ebenso werden die aus verschiedenen Kombinationen gewonnenen Bestimmungen eine ungleiche Genauigkeit besitzen.

Wären im zweiten oder dritten Fall die Funktionen \(V,V',V''\) etc. außerdem so beschaffen, dass \(\pi-\rho+1\) von ihnen oder mehr als Funktionen der übrigen angesehen werden könnten, so wäre die Aufgabe in Bezug auf die letzteren Funktionen noch immer mehr als bestimmt, in Bezug auf die Unbekannten \(x,y,z\) etc. dagegen unbestimmt; deren Werte ließen sich nämlich selbst dann nicht bestimmen, wenn die Werte der Funktionen \(V,V',V''\) etc. absolut genau gegeben wären: Diesen Fall schließen wir aber von unserer Untersuchung aus.

Sind \(V,V',V''\) etc. an sich keine \textit{linearen} Funktionen ihrer Unbestimmten, so wird dies erreicht, wenn anstelle der ursprünglichen Unbekannten ihre Differenzen von Näherungswerten eingeführt werden. Diese dürfen als anderweitig bekannt angenommen werden. Die bei den Bestimmungen \(V=L,V'=L',V''=L''\) etc. zu befürchtenden mittleren Fehler bezeichnen wir jeweils mit \(m,m',m''\) etc., die Gewichte der Bestimmungen mit \(p,p',p''\) etc., sodass \(pmm=p'm'm'=p''m''m''\) etc. gilt. Wir setzen das Verhältnis der mittleren Fehler zueinander als bekannt voraus, sodass die Gewichte, von denen eines nach Belieben angenommen werden kann, bekannt sind. Schließlich setzen wir
\SourceLeaf{027}
\[
\begin{aligned}
(V-L)\BareRoot{p}&=v,\quad(V'-L')\BareRoot{p'}=v',\\
(V''-L'')\BareRoot{p''}&=v''\text{ etc.}
\end{aligned}
\]
Die Sache verhält sich daher offensichtlich so, als hätten unmittelbare Beobachtungen gleicher Genauigkeit, nämlich mit dem mittleren Fehler \(=m\BareRoot{p}=m'\BareRoot{p'}=m''\BareRoot{p''}\) etc. oder mit dem zugeteilten Gewicht \(=1\), die Werte
\[
v=0,\quad v'=0,\quad v''=0\text{ etc.}
\]
geliefert.

\clearpage
